\chapter{Context: Domain and Literature}



% title: Fulfilling Information Needs in Knowledge-Intensive Domains: Entity-Oriented Strategies for Information Extraction and Access

% concepts of information need and information object
% information extraction and access
% bit of context of nlp, ir, semantic web and databases
% description of the domain and subdomains (law judgements and data from criminal investigations)
% literature of most relevant tasks (NER, NEL, coreference resolution, QA) 
\todo{who access the information? humans, llms, humans through llms accessing}
This chapter lays the foundations for the rest of the thesis by carefully introducing the background on which our work builds. We begin by discussing about information access and the notions of \emph{information needs} and \emph{information objects}~\cite{entityorientedsearch2018balog}. We define what is an entity and then turn to information extraction. To better situate the discussion within the relevant branches of computer science---namely \acl{nlp}, \acl{ir}, \acl{sw}, and \acl{db}---a brief historical overview is presented for each of them, focusing on the interceptions with this work\todo{improve}. Once this broader context is established, we describe the specific knowledge-intensive domains that motivate our study, namely legal judgments and data from criminal investigations. Finally, we review the literature on the concrete tasks that are most directly related to our research---such as \acl{er}, \acl{el}, \acl{cr}, and \acl{qa}.

Although part of this \phd work has been applied to general-domain benchmark data, the motivaton behind still resides in the peculiarities of specific domains and their properties.\todo{fix and enumerate the properties!}
For this reason, concepts and tasks definitions from the literature are paired with examples referring to the focused domains with the twofold aim of facilitating the reader comprehension of the concepts and familiarization with the domains.  

\todo{guide the reader through info extraction/text annotation/ner, coref,... add image of the taxonomy?}
\todo{rethink the "in this section part"}



